% =====================================================================
%  APPENDIX: The charge circle -- winding equals ladder number
%  ADDED 2026-07-12. Discharges the first lemma of
%  op:dictionary-lemmas (sec:algebraic-dictionary) in the same revision
%  that posed it. Proof rests only on the composition-algebra identity
%  and the standard Clifford rotor fact; verified to machine precision
%  in a convention-independent implementation (Cayley--Dickson table;
%  the third plane's (1,7,6) orientation emerged with the correct
%  sign). CUTOVER: the verification is slated as notebook P of the
%  notebook record; add \ref{app:nb-P} here when the record is input.
% =====================================================================

\section{The charge circle: winding equals ladder number}
\label{app:winding-number}

This appendix proves the charge row of the dictionary of \S\ref{sec:algebraic-dictionary}:
Furey's electric-charge operator --- one third of the fermionic ladder count --- and the
fibration's winding number are one operator in two coordinatisations. The proof is four steps,
and every step rests only on structure this paper has already paid for.

\emph{Setup.} Let $\ell$ be the selected imaginary unit ($e_{1}$ in the conventions of
Appendix~\ref{app:setup}) and $J = L_{\ell}$ the complex structure it induces, splitting
$\mathbb{O} = \mathbb{C}\ell \oplus \mathbb{C}^{3}$ with the three colour planes the Fano lines
through $\ell$: $(e_{2}, e_{3})$, $(e_{4}, e_{5})$, $(e_{7}, e_{6})$, each oriented by
$y_{k} = J x_{k}$. Write $c_{j} = L_{e_{j}}$ for the six left multiplications of the complement.
The composition identity gives, in one line, the Clifford relations: $L_{a}^{2} = -|a|^{2}$ for
imaginary $a$ by alternativity, and polarising, $\{c_{a}, c_{b}\} = -2\langle a, b\rangle$. The
six $c_{j}$ therefore generate $\mathbb{C}l(6)$ acting irreducibly on
$\mathbb{C} \otimes \mathbb{O} \cong \mathbb{C}^{8}$ --- Furey's arena, reached here from the
fibration side.

\emph{Modes.} For each colour plane set
$\alpha_{k} = \tfrac12\,(c_{x_{k}} - \mathrm{i}\, c_{y_{k}})$. The Clifford relations give
$\alpha_{k}^{2} = 0$ and $\{\alpha_{j}, \alpha_{k}^{\dagger}\} = \delta_{jk}$: three fermionic
modes, one per colour plane. Their Fock grading of $\mathbb{C}^{8}$ is
$\Lambda^{0} \oplus \Lambda^{1} \oplus \Lambda^{2} \oplus \Lambda^{3} = 1 + 3 + 3 + 1$ --- the
multiplet row of the dictionary, and the grading of Open Problem~\ref{op:ideal-orbit}. Write
$N = \sum_{k} \alpha_{k}^{\dagger}\alpha_{k}$ for the total number operator; Furey's electric
charge is $Q = N/3$ on the particle ideal.

\emph{The engine.} A two-line computation from the relations:
\begin{equation}
\alpha_{k}^{\dagger}\alpha_{k}
 = -\tfrac14\,(c_{x_{k}} + \mathrm{i} c_{y_{k}})(c_{x_{k}} - \mathrm{i} c_{y_{k}})
 = \tfrac12\,\bigl(1 + \mathrm{i}\, c_{x_{k}} c_{y_{k}}\bigr),
\qquad\text{i.e.}\qquad
\mathrm{i}\, c_{x_{k}} c_{y_{k}} = 2 N_{k} - 1 .
\label{eq:winding-engine}
\end{equation}
The number operator of each mode \emph{is} the Clifford bilinear of its plane.

\emph{The lift.} The rotation by angle $\theta$ in the plane $(x_{k}, y_{k})$ acts on the
module by the rotor $\exp(\tfrac{\theta}{2}\, c_{x_{k}} c_{y_{k}})$ --- conjugation by it sends
$c_{x_{k}} \mapsto \cos\theta\, c_{x_{k}} + \sin\theta\, c_{y_{k}}$, the standard half-angle
fact. The \emph{colour circle} --- the determinant $U(1)$ of the $J$-commuting rotations
$U(3)$, each colour plane turned once with the record plane fixed --- therefore acts on the
ideal as
\begin{equation}
\exp\Bigl(\tfrac{\theta}{2} \sum_{k} c_{x_{k}} c_{y_{k}}\Bigr)
 = \exp\bigl(-\mathrm{i}\,\theta\,(N - \tfrac32)\bigr),
\label{eq:winding-lift}
\end{equation}
by~\eqref{eq:winding-engine}. The right side is diagonal in the Fock basis with winding numbers
$W = \tfrac32 - N$ per full turn.

\begin{lemma}[Winding equals ladder number]\label{lem:winding-number}
On the ideal, the ladder number operator is the relative winding of the colour circle's spinor
lift, $N = W_{\mathrm{vac}} - W$ with $W_{\mathrm{vac}} = \tfrac32$ the vacuum's winding, so that
\begin{equation}
Q \;=\; \frac{N}{3} \;=\; \frac{W_{\mathrm{vac}} - W}{3}.
\end{equation}
Furey's charge operator and the fibration's winding are one operator in two coordinatisations:
the charge thirds are winding integers divided by the number of colour planes, and the ladder's
integer spacing is the winding's.
\end{lemma}

Three remarks. First, the half-integer vacuum winding is not a blemish but the spinor double
cover announcing itself --- the same $\mathbb{Z}_{2}$ that makes spin one half --- and it drops
out of every physical (relative) statement. Second, adjoining the record plane's rotation ---
the full $J$-circle in place of the colour circle --- contributes the mode-zero bilinear, an
\emph{integer} unit shift: the leptonic sector's integer charges, whose alignment with the
register normalisation of \S\ref{sec:charges} is the dictionary's next row. Third, on the
conjugate ideal the windings negate and the antiparticle charges follow. The proof above uses
only the composition identity and the rotor fact; it has been verified to machine precision in
a convention-independent implementation, the multiplication table generated by Cayley--Dickson
doubling and the three planes selected by $J$ itself.
